\section{Separation, quadratic mass, and angular energy}
\label{sec:current-energy}

The two separation arguments below produce a positive current, possibly with
an anti-invariant correction.  We first control its trace measure at the
two-dimensional scale.  An intersection pairing then controls the correction
and the variation of the complex-line directions.  These estimates will
permit the density decomposition in \Cref{sec:density-splitting}.

Separation by positive currents follows the framework of Sullivan and
Harvey--Lawson \cite{Sullivan1976,HarveyLawson1983}, with the
almost-complex formulation of Li--Zhang
\cite[Section~3]{LiZhang2009}. The closed lift,
quadratic mass estimate and resolvent argument also occur in
\emph{Taming implies compatibility on four-manifolds}
\cite[Sections~2--4]{OpenAITaming2026}. We develop the estimates for
both separators together. In the prescribed-class application, the
closed current can have nonzero cohomology class, and its square must
remain in the energy identity.

\subsection{Conventions and separation}

Fix a smooth almost-complex structure $J$ on the closed four-manifold $X$,
give $X$ its complex orientation, and choose a smooth $J$-Hermitian metric
$g$.  Its fundamental form is
\[
 G(u,v)=g(Ju,v).
\]
The pointwise orthogonal decomposition is
\begin{equation}\label{eq:current-bundles}
 \Lambda^2=I\oplus E,
 \qquad I=\Lambda_J^+=\R G\oplus\Lambda_g^-,
 \qquad E=\Lambda_J^-\subset\Lambda_g^+.
\end{equation}
Let $R$ denote the orthogonal projection onto $E$.  A unit complex-line
bivector is $\ell=e\wedge Je$, where $|e|_g=1$.  We use $g$ to identify
bivectors and two-forms.  Consequently
\begin{equation}\label{eq:current-line-normalization}
 |\ell|_g=1,\qquad \ell^+=G/2,\qquad G(\ell)=1,
 \qquad |G|_g^2=2.
\end{equation}
The bundle of these bivectors is denoted by $\mathcal L\longrightarrow X$.

A two-current is a continuous real linear functional on
$\Omega^2(X)=C^\infty(X,\Lambda^2)$.  Under the metric identification its
coefficient is a distributional two-form, denoted by the same letter:
$T(\xi)=\langle T,\xi\rangle$.  Thus
\begin{equation}\label{eq:current-dual-convention}
 T\text{ is closed}\ \Longleftrightarrow\
 T(\dd\alpha)=0\text{ for all }\alpha\in\Omega^1(X)
 \ \Longleftrightarrow\ \dd^*T=0
 \ \Longleftrightarrow\ \dd(*T)=0.
\end{equation}
The cohomology class dual to a closed current is $[*T]\in H^2(X;\R)$.
For smooth coefficients this convention gives
$T(\xi)=\int_X\xi\wedge *T$.  Inner products between $L^2$ coefficients
will be written $\langle\ ,\ \rangle_2$.

Let
\[
 \mathcal P_J=\{\xi\in\Omega^2(X):\xi(\ell)>0
                    \text{ for every }\ell\in\mathcal L\},
 \qquad
 \mathcal Z_J=\{\xi\in C^\infty(X,I):\dd\xi=0\}.
\]
The set $\mathcal P_J$ is an open convex cone in the smooth topology;
openness already holds in the uniform topology.  Its closed elements are
precisely the symplectic forms taming $J$.  The elements of
$\mathcal P_J\cap\mathcal Z_J$ are the compatible symplectic forms.

\begin{lemma}[The two separators]\label{lem:current-separators}
The following statements hold.
\begin{enumerate}[label=\textup{(\roman*)}]
\item If $\mathcal P_J\cap\mathcal Z_J$ is empty, there is a nonzero current
 $N$ which annihilates $\mathcal Z_J$ and all anti-invariant forms, and is
 nonnegative on semipositive invariant forms.
\item If $H\in H^2(X;\R)$ has no representative in $\mathcal P_J$, there
 is a nonzero positive invariant current $N$ which is closed and satisfies
 $N(H)\le0$.
\end{enumerate}
In either case one can normalize $N(G)=1$ and choose a positive Borel
measure $\lambda$ on $\mathcal L$ such that
\begin{equation}\label{eq:current-line-measure}
 N(\xi)=\int_{\mathcal L}\xi_y(\ell)\,\dd\lambda(y,\ell),
 \qquad \mu=(\operatorname{pr}_X)_*\lambda,
 \qquad \mu(X)=1.
\end{equation}
In particular $N$ has order zero and belongs to $H^{-3}(X,\Lambda^2)$.
\end{lemma}

\begin{proof}
For \textup{(i)}, apply separation of an open convex set and a disjoint
linear subspace in the real locally convex space $C^\infty(X,I)$.
The separating continuous functional vanishes on $\mathcal Z_J$ and is
nonnegative on its positive cone.  Extend it to all two-forms by
precomposing with $1-R$.  Adding a positive multiple of $G$ and taking
a limit proves nonnegativity on every semipositive invariant form.

For \textup{(ii)}, choose a smooth closed representative $\beta$ of $H$
and separate the open cone $\mathcal P_J$ from the affine subspace
$\beta+\dd\Omega^1(X)$ in $\Omega^2(X)$.  If $N$ is the separating
functional, the unbounded translation directions in this affine subspace
give $N(\dd\alpha)=0$.  Every positive form can be multiplied by any
positive scalar, so the separation inequalities imply both
$N(\mathcal P_J)\ge0$ and $N(\beta)\le0$.  Adding an arbitrary real
multiple of an anti-invariant form to an element of $\mathcal P_J$ shows
that $N$ annihilates anti-invariant forms.  The same limiting argument
gives positivity on semipositive invariant forms.

For either functional, comparison of an invariant form with sufficiently
large positive multiples of $G$ gives
\[
 |N(\xi)|\le C N(G)\|\xi\|_{C^0}.
\]
If $N(G)$ were zero, this inequality and anti-invariance would make the
functional zero.  Normalize $N(G)=1$.  Locally, the functional is therefore
represented by a positive Hermitian matrix of measures.  Its trace is a
positive measure $\mu$; the Radon--Nikodym matrix relative to $\mu$ is
positive semidefinite with trace one.  A measurable spectral decomposition
of this $2\times2$ matrix expresses it as a convex combination of rank-one
Hermitian projections.  Borel local frames and a Borel partition of $X$
give the global measure $\lambda$ in
\Cref{eq:current-line-measure}.  Finally $H^3\hookrightarrow C^0$ in real
dimension four, so every such measure coefficient belongs to $H^{-3}$.
\end{proof}

\subsection{A closed lift of anti-invariant forms}

\begin{lemma}[Closed lift]\label{lem:closed-lift}
There is a classical pseudodifferential operator
$D:C^\infty(X,E)\longrightarrow\Omega^2(X)$ of order zero with
\begin{equation}\label{eq:current-lift-identities}
 \dd D=0,\qquad RD=\Id.
\end{equation}
It maps distributions to distributions.  For the separator in
\Cref{lem:current-separators}\textup{(i)}, define
\begin{equation}\label{eq:current-correction}
 Q(\xi)=-N(DR\xi),\qquad T=N+Q.
\end{equation}
Then $Q$ is anti-invariant and self-dual, $Q,T\in H^{-3}$, and $T$
annihilates every closed smooth two-form.  In particular $T$ is a closed
current and $[*T]=0$.

For a smooth family of $(J,g)$ with constant dimension of the space of
closed anti-invariant forms, these lifts can be chosen smoothly in the
parameter as operators between the corresponding smooth bundles.
\end{lemma}

\begin{proof}
On sections of $E$ consider $P=R\dd\dd^*$.  If $a$ is self-dual, then
the self-dual projection of $\xi\wedge\iota_\xi a$ is
$|\xi|^2a/2$: in a frame with first covector $\xi/|\xi|$, the two
self-dual halves of $a$ have equal length.  Hence
\[
 \sigma_2(P)(x,\xi)=\tfrac12|\xi|_g^2\Id_E,
 \qquad \langle Pa,a\rangle_2=\|\dd^*a\|_2^2.
\]
Thus $P$ is elliptic, self-adjoint, and nonnegative.  Its kernel consists
of smooth sections with $\dd^*a=0$.  Since $a=*a$, these sections also
satisfy $\dd a=0$, and hence are harmonic.  Conversely a closed
anti-invariant form is self-dual and therefore harmonic and in the kernel.

Let $H_0$ be the orthogonal projection onto this finite-dimensional
kernel, and let $L$ be the inverse of $P$ on its orthogonal complement,
extended by zero on the kernel. Elliptic Fredholm theory and the
parametrix give $L$ as a classical operator of order $-2$, with
$H_0$ smoothing and $PL=\Id-H_0$; see
\cite[Sections~4--5 and~10]{TaylorPseudodifferentialNotes}.
Indeed, if $B_0P=\Id+K$ is a parametrix identity with $K$ smoothing,
then $B_0(\Id-H_0)=L+KL$. Elliptic regularity makes $KL$ smoothing,
which gives the asserted order of $L$. Set
\begin{equation}\label{eq:current-lift-definition}
 D=\dd\dd^*L+H_0.
\end{equation}
The image of both summands is closed, and
$RD=PL+H_0=\Id$.  This proves
\Cref{eq:current-lift-identities}.

For closed $\xi$, the form $\xi-DR\xi$ is closed and invariant, so
$T(\xi)=N(\xi-DR\xi)=0$.  The functional $Q$ only depends on $R\xi$,
which proves its anti-invariance and self-duality.  Its $H^{-3}$ membership
follows because an order-zero operator and its adjoint act boundedly on
every Sobolev space
\cite[Proposition~5.5]{TaylorPseudodifferentialNotes}.
The vanishing on exact forms makes $T$ closed.
Its pairings with all closed two-forms vanish, so Poincare duality, or the
harmonic projection of $*T$, gives $[*T]=0$.

For the family assertion, identify the varying bundles locally in the
parameter and use a smooth unitary identification of the $L^2$ spaces.
The elliptic operators then have common domain $H^2$ in a fixed
Hilbert space. Choose a spectral contour enclosing only zero at one
parameter value. The resolvent identity and elliptic regularity give
a smooth family of resolvents on this contour after shrinking the
parameter neighborhood. Its spectral projection has locally constant
rank. Since the kernel dimension is constant with that rank, the
enclosed spectral subspace consists exactly of the kernel. Thus $H_0$
depends smoothly on the parameter, and so does
\[
 L=(P+H_0)^{-1}(\Id-H_0)
\]
on every Sobolev scale. \Cref{eq:current-lift-definition} gives the
required smooth lifts.
\end{proof}

The projected elliptic operator also underlies Lejmi's local
compatibility construction \cite[Theorem~2.1]{Lejmi2006}; the related
global linear correction appears in
\cite[Proposition~3.1]{DraghiciZhang2012}. The harmonic projection in
\Cref{eq:current-lift-definition} records the possibly nonzero kernel.
The lift provides closedness and the prescribed anti-invariant part.
The estimates below will address existence of a positive closed
invariant form.

From now on we use one of the two following setups:
\begin{enumerate}[label=\textup{(\Alph*)}]
\item $N$ is the normalized separator in
 \Cref{lem:current-separators}\textup{(i)}, and $Q,T$ are given by
 \Cref{eq:current-correction};
\item $N$ is a normalized positive invariant closed current, $Q=0$, and
 $T=N$.
\end{enumerate}
The second setup includes \Cref{lem:current-separators}\textup{(ii)}.
In both cases $T$ is closed, but the obstructions to be removed are
different. In setup \textup{(A)}, we will prove $Q=0$: the positive
current $N=T$ would then annihilate every closed form, which is
impossible when $J$ has a closed tamer. In setup \textup{(B)}, a class separator
satisfies $N(H)\le0$; the task is instead to prove $N(H)>0$ from the
hypotheses on the proposed taming class. The estimates below apply to
both setups. They first put $Q$ in $L^2$ and control nearby line
directions; \Cref{sec:density-splitting} will then eliminate the two
obstructions.

Constants below can depend on $J,g$ and
the fixed operators, but are independent of the center, the small scale,
and the choice of submeasure of $\lambda$.

\subsection{The quadratic mass estimate}

We will use the following elementary consequence of the order-zero
calculus.  The proof records the dependence on the radial scale.

\begin{lemma}[Radial order-zero estimate]\label{lem:current-radial-operator}
Let $A$ be a fixed classical operator of order zero between smooth vector
bundles on a compact four-manifold.  Fix a coordinate radius below the
injectivity radius.  Suppose that smooth sections $a_{s,p}$, supported in
that coordinate neighborhood of $p$, satisfy, for $0<s<s_0$,
\[
 |\nabla^j a_{s,p}(y)|\le C_j(s+\dist(p,y))^{-a-j},
 \qquad 0\le j\le6,
 \qquad a\in\{0,1\},
\]
with uniform boundedness on a fixed outer annulus.  Then, for
$\delta=s+\dist(p,x)$ small,
\begin{equation}\label{eq:current-radial-operator}
 |Aa_{s,p}(x)|\le
 \begin{cases}
 C\delta^{-1},&a=1,\\
 C(1+|\log\delta|),&a=0.
 \end{cases}
\end{equation}
Away from a smaller fixed neighborhood of $p$, the output is uniformly
bounded in $s$ and $p$.
\end{lemma}

\begin{proof}
We use two precise facts of the order-zero calculus: in local left
quantization its symbol is bounded uniformly in the covariable, and its
off-diagonal kernel satisfies
$|K_A(x,y)|\le C\dist(x,y)^{-4}$ for small positive distance.  Smooth
remainders have uniformly bounded kernels.  These statements hold for
bundle-valued classical operators, after finitely many coordinate
localizations; see
\cite[Section~2, Propositions~2.1--2.2]{TaylorPseudodifferentialNotes}.

For a fixed evaluation point $x$, split the input using a smooth cutoff
supported in $B_{c\delta}(x)$ and equal to one on
$B_{c\delta/2}(x)$, with a fixed small $c$.  On its support
$s+\dist(p,y)$ is comparable to $\delta$.  After translation and dilation
by $\delta$, the localized input has derivatives through order six
bounded by $C\delta^{-a}$.  Integration by parts in its Fourier transform
therefore bounds its Fourier $L^1$ norm by $C\delta^{-a}$, since
$(1+|\xi|)^{-6}$ is integrable in dimension four.  The bounded symbol
gives the same bound for the local contribution to $Aa_{s,p}(x)$.

The remaining part has distance at least $c\delta/2$ from $x$.
Its portion in $B_{C\delta}(p)$ contributes at most
\[
 C\delta^{-4}\int_{B_{C\delta}(p)}
          (s+\dist(p,y))^{-a}\,\dd V_g(y)
 \le C\delta^{-a}.
\]
On a shell of radius $r=2^j\delta\ge C\delta$ about $x$, distance from
$p$ is comparable to $r$, so its contribution is at most
$Cr^{-4}r^4r^{-a}=Cr^{-a}$.  The shells up to the fixed coordinate
radius sum to $C\delta^{-1}$ if $a=1$, and to
$C(1+|\log\delta|)$ if $a=0$.  The rest is uniformly bounded.
If $x$ stays away from $p$, the singular input near $p$ has uniformly
bounded $L^1$ norm and the kernel there is bounded; the remaining input
is uniformly smooth.  This proves the final assertion.
\end{proof}

\begin{lemma}[Quadratic trace mass]\label{lem:quadratic-mass}
In either setup \textup{(A)} or \textup{(B)}, there is a constant $C$
such that
\begin{equation}\label{eq:current-quadratic-mass}
 \mu(B_s(p))\le Cs^2\qquad(p\in X,\ 0<s\le1).
\end{equation}
\end{lemma}

\begin{proof}
Fix first a radius $\rho_0>0$ for uniform normal coordinates and a smooth
radial cutoff which is one on $B_{2\rho_0}(p)$ and supported in
$B_{3\rho_0}(p)$.  This cutoff will not change when a smaller inner
radius is chosen below. The logarithmic potential will detect the mass
of $B_s(p)$, while a square-root potential will dominate the
inverse-distance errors caused by varying $J$ and by the closed lift.
With $t=\dist(p,\cdot)$ define the globally
smooth cut-off functions
\[
 u_s=\chi(t)\log\sqrt{s^2+t^2},\qquad
 h_s=\chi(t)\sqrt{s^2+t^2},\qquad
 \dd_J^cu=-\dd u\circ J.
\]
Smoothness at $p$ follows from the dependence on $t^2$ and from $s>0$.
For a function $f$, put
\[
 \beta(f)=
 \begin{cases}
  (1-DR)\dd\dd_J^cf,&\text{in setup \textup{(A)}},\\
  \dd\dd_J^cf,&\text{in setup \textup{(B)}}.
 \end{cases}
\]
In setup \textup{(A)}, $\beta(f)$ is closed and invariant, and in setup
\textup{(B)} it is exact.  Hence $N(\beta(f))=0$ in both cases.

The principal symbol of $\dd\dd_J^c$ on functions is a multiple of
$\xi\wedge J^*\xi$, which is $J$-invariant.  Thus
$R\dd\dd_J^c$ is a first-order operator.  Differentiating the radial
functions gives, for $0\le j\le6$,
\[
 |\nabla^jR\dd\dd_J^cu_s|\le C_j(s+t)^{-1-j},\qquad
 |\nabla^jR\dd\dd_J^ch_s|\le C_j(s+t)^{-j}.
\]
On the cutoff annulus these inputs and their derivatives are uniformly
bounded.  \Cref{lem:current-radial-operator} applied to $D$ therefore
bounds the two correction terms by
$C/(s+t)$ and $C(1+|\log(s+t)|)$ respectively.

For clarity, the positive part of the test can be computed in a fixed
Hermitian Euclidean space.  If $f=f(t)$ and $\ell$ is a unit complex
line, its Hessian trace is
\[
 \dd\dd_{J_p}^cf(\ell)
 =2f'(t)/t+\bigl(f''(t)-f'(t)/t\bigr)a,
 \qquad 0\le a\le1,
\]
where $a$ is the squared length of the radial unit vector projected onto
the line.  For the two functions without cutoff this gives
\begin{align}
 \dd\dd_{J_p}^c\log\sqrt{s^2+t^2}(\ell)
   &\ge \frac{2s^2}{(s^2+t^2)^2},\label{eq:current-log-trace}\\
 \dd\dd_{J_p}^c\sqrt{s^2+t^2}(\ell)
   &\ge \frac1{\sqrt{s^2+t^2}}.\label{eq:current-root-trace}
\end{align}
In normal coordinates $g-g_p=O(t^2)$ and $J-J_p=O(t)$, while the first
derivatives of $J$ are bounded.  Replacing the frozen complex line and
operator by the actual ones changes these estimates by at most
$C/(s+t)$ for $u_s$ and $C$ for $h_s$.  Indeed the Hessian errors are
bounded by $Ct|\nabla^2f|$ and the coefficient-derivative errors by
$C|\nabla f|$.  With $\delta=s+t$, it follows that on
$B_{\rho_0}(p)$
\begin{align}
 \beta(u_s)(\ell)&\ge c_1s^2\delta^{-4}-C_1\delta^{-1},
                  \label{eq:current-log-corrected}\\
 \beta(h_s)(\ell)&\ge c_2\delta^{-1}
                   -C_2(1+|\log\delta|).
                  \label{eq:current-root-corrected}
\end{align}
The inequalities include setup \textup{(B)}, where the correction terms
are absent.

All constants just obtained use the already fixed cutoff.  Now choose
$0<\rho<\rho_0$ so small that
\[
 C_2(1+|\log\delta|)\le c_2/(2\delta)
                  \qquad(0<\delta\le2\rho).
\]
Next choose a fixed multiplier $M\ge2C_1/c_2$.  If $s\le\rho$ and
$t\le\rho$, \Cref{eq:current-log-corrected,eq:current-root-corrected}
give $\beta(u_s+Mh_s)(\ell)\ge c_1s^2\delta^{-4}$.
On $X\setminus B_\rho(p)$ all terms are bounded independently of $s$,
including the nonlocal corrections, by the last assertion of
\Cref{lem:current-radial-operator}.  Since $\delta\le2s$ on $B_s(p)$,
we conclude globally that
\[
 \beta(u_s+Mh_s)(\ell)\ge c s^{-2}\mathbf1_{B_s(p)}-C.
\]
Apply $N$ and use $N(\beta(u_s+Mh_s))=0$ and $\mu(X)=1$ to obtain
\Cref{eq:current-quadratic-mass} for $s\le\rho$.  For larger $s$ the
same assertion follows from $\mu(X)=1$.  Compactness makes all choices
uniform in $p$.
\end{proof}

\subsection{The resolvent kernel and positive pairings}

We now have a uniform mass bound at the two-dimensional scale. The
next step is to compare directions of nearby complex lines. Smoothing
allows this comparison through an intersection pairing while keeping
the error controlled by the mass bound just proved.

Let $\Delta=\dd\dd^*+\dd^*\dd$ be the nonnegative Hodge Laplacian and
set
\begin{equation}\label{eq:current-smoothing}
 S_s=(1+s^2\Delta)^{-3},\qquad 0<s\le s_0.
\end{equation}
It commutes with $\dd,\dd^*$ and $*$.  At each fixed scale it maps
$H^{-3}$ to $H^3$, which suffices for all the pairings below.  The kernel
of $S_s^2$ is continuous, being of order $-12$ in dimension four.

\begin{lemma}[Resolvent kernel]\label{lem:current-kernel}
Choose a fixed radius $r_0$ smaller than the injectivity radius, and let
$\tau_{z\to y}$ be parallel transport along the minimizing geodesic
when $\dist(y,z)<r_0$.  The kernel $K_s$ of $S_s^2$ can be written
\begin{equation}\label{eq:current-kernel-split}
 K_s(y,z)=k_s(y,z)\tau_{z\to y}+E_s(y,z),
\end{equation}
where $k_s$ is nonnegative and supported where $\dist(y,z)<r_0$.
For each integer $k\ge0$ there are constants such that
\begin{align}
 k_s(y,z)&\ge cs^{-4}&&\text{if }\dist(y,z)<s,
                         \label{eq:current-kernel-lower}\\
 k_s(y,z)&\le C_ks^{-4}(1+\dist(y,z)/s)^{-k},
                         \label{eq:current-kernel-upper}\\
 |E_s(y,z)|&\le C_ks^{-2}(1+\dist(y,z)/s)^{-k}.
                         \label{eq:current-kernel-error}
\end{align}
The small upper bound on $s_0$ is fixed independently of $y,z$.
\end{lemma}

\begin{proof}
The spectral theorem gives the operator identity
\begin{equation}\label{eq:current-gamma-resolvent}
 S_s^2=\frac1{5!}\int_0^\infty h^5e^{-h}e^{-s^2h\Delta}\,\dd h.
\end{equation}
We recall the precise heat-kernel facts being used for the nonnegative
Hodge Laplacian on the closed manifold. Its local small-time kernel
has a parametrix
\[
 \chi(d)(4\pi t)^{-2}e^{-d^2/(4t)}
       \bigl(a_0(y,z)+ta_1(y,z)+\cdots+t^ma_m(y,z)\bigr),
 \qquad d=\dist(y,z),
\]
with arbitrarily high-order uniform remainder after increasing $m$;
see \cite[Theorems~1.1 and~3.1]{Ludewig2019}.
Here $\chi$ is supported inside the injectivity radius and equals one
near the diagonal.  For the Hodge Laplacian
$a_0(y,z)=j(y,z)^{-1/2}\tau_{z\to y}$, with $j$ smooth, strictly
positive, and equal to one on the diagonal
\cite[Section~3, equation~(3.4), author version]{Ludewig2018}.
After subtracting the leading
term, for every prescribed $L$ the kernel is bounded, for
$0<t<t_0$, by
\[
 Ct^{-1}e^{-d^2/(Ct)}+C_Lt^L.
\]
The local expansion gives the Gaussian part of this bound; away from
the diagonal, the global Gaussian estimate makes the cutoff error
smaller than every power of $t$
\cite[Theorem~3.5]{Ludewig2019}. For $t\ge t_0$, the heat kernel is
uniformly bounded: factor the heat operator through two fixed
positive-time smoothing operators and the $L^2$ contraction
$e^{-(t-t_0)\Delta}$. This last assertion uses the nonnegativity and
self-adjointness of the Hodge Laplacian.

Take $\chi\ge0$ and use its leading coefficient to define
\[
 k_s(y,z)=\frac{\chi(d)j(y,z)^{-1/2}}{5!(4\pi)^2}s^{-4}
       \int_0^\infty h^3e^{-h-d^2/(4s^2h)}\,\dd h.
\]
For $d<s$, integration over $1\le h\le2$ gives
\Cref{eq:current-kernel-lower}.  For every fixed $k$, the elementary
inequality
$e^{-r^2/(ch)}\le C_k(1+r)^{-k}(1+h^{k/2})$ shows that the last
integral is bounded by $C_k(1+d/s)^{-k}$.  This proves
\Cref{eq:current-kernel-upper}.

Integrating the $t^{-1}$ heat-kernel error in
\Cref{eq:current-gamma-resolvent} gives
\[
 Cs^{-2}\int_0^\infty h^4e^{-h-d^2/(Cs^2h)}\,\dd h
 \le C_ks^{-2}(1+d/s)^{-k}.
\]
For the uniform remainder, integration over $s^2h<t_0$ gives
$C_Ls^{2L}$.  Since $d\le\operatorname{diam}X$, this is bounded by
the right side of \Cref{eq:current-kernel-error} when
$2L\ge k-2$.  On $s^2h\ge t_0$, the factor $e^{-h}$ makes both the
heat-kernel contribution and the extended leading term smaller than
every power of $s$.  This proves the claimed error for each $k$.
\end{proof}

For any positive submeasure $\lambda_i\le\lambda$, let $N_i$ be its
current and $\mu_i$ its projection to $X$.  These currents need not be
closed.  Define
\begin{equation}\label{eq:current-error-bracket}
 B_{s,\mu_i}(z)=s^{-2}\int_X
       (1+\dist(y,z)/s)^{-6}\,\dd\mu_i(y),
 \qquad B_s=B_{s,\mu}.
\end{equation}
By \Cref{lem:quadratic-mass}, a dyadic-shell sum gives
\begin{equation}\label{eq:current-bracket-bound}
 0\le B_{s,\mu_i}(z)\le B_s(z)\le C.
\end{equation}
Indeed the shell of radius $2^js$ contributes at most $C2^{-4j}$;
the estimate for radii above a fixed coordinate radius follows from
$\mu(X)=1$.

\begin{lemma}[Positive-current pairing]\label{lem:current-positive-pairing}
For any $\lambda_1,\lambda_2\le\lambda$ and their currents,
\begin{align}
 \langle S_sN_1,*S_sN_2\rangle_2
 &\ge cs^{-4}
       \iint_{\dist(y,z)<s}|\ell-\tau_{z\to y}m|^2
             \,\dd\lambda_1(y,\ell)\dd\lambda_2(z,m)
       \notag\\[-2pt]
 &\hspace{8mm}-C\int_X B_{s,\mu_1}(z)\,\dd\mu_2(z),
                                  \label{eq:current-positive-pairing}\\
 \|S_sN_i\|_2&\le C/s.            \label{eq:current-submeasure-ltwo}
\end{align}
One can replace $B_{s,\mu_1}$ in the lower bound by $B_s$.
\end{lemma}

\begin{proof}
Parallel transport preserves the Hodge star.  From
\Cref{eq:current-line-normalization}, direct expansion gives
\begin{equation}\label{eq:current-angular-identity}
 \langle\ell,*\tau_{z\to y}m\rangle
 =\tfrac12|\ell-\tau_{z\to y}m|^2
       -\tfrac14|G_y-\tau_{z\to y}G_z|^2.
\end{equation}
For example, the right side equals
$-\langle\ell,\tau m\rangle+\tfrac12\langle G_y,\tau G_z\rangle$,
which is the left side after splitting into self-dual and anti-self-dual
parts.  The last squared difference is bounded by $Cd^2$, where
$d=\dist(y,z)$.

Express the pairing using the continuous kernel of $S_s^2$ and
\Cref{eq:current-line-measure}.  Its leading scalar is nonnegative, so
\Cref{eq:current-kernel-lower,eq:current-angular-identity} supply the
positive term.  The $Cd^2$ term, using
\Cref{eq:current-kernel-upper} with exponent at least eight, and the
error kernel, using \Cref{eq:current-kernel-error} with exponent six,
are bounded in absolute value by
\[
 Cs^{-2}\iint(1+d/s)^{-6}\,\dd\mu_1(y)\dd\mu_2(z).
\]
This is the asserted error.  The norm estimate follows by using the
absolute kernel bound in the unstarred pairing:
\[
 \|S_sN_i\|_2^2
 \le Cs^{-4}\iint(1+d/s)^{-6}\,\dd\mu_i(y)\dd\mu_i(z)
 \le Cs^{-2}\mu_i(X)\le Cs^{-2}.
\]
The kernel formulas are legitimate for measure coefficients because the
kernel is continuous at fixed $s$; alternatively they follow by smooth
approximation at that fixed scale.
\end{proof}

\subsection{Projection leakage and the topological energy}

\begin{lemma}[Projection commutator]\label{lem:current-commutator}
If $\Pi$ is a smooth orthogonal bundle projection and $\Pi W=0$ for a
distributional form $W$ with $S_sW\in L^2$, then
\begin{equation}\label{eq:current-projection-leakage}
 \|\Pi S_sW\|_2\le Cs\|S_sW\|_2.
\end{equation}
Consequently, for every positive submeasure current $N'\le_\lambda N$
and every anti-invariant distribution $Z$ with $S_sZ\in L^2$,
\begin{equation}\label{eq:current-mixed-bound}
 |\langle S_sN',S_sZ\rangle_2|\le C\|S_sZ\|_2.
\end{equation}
Here $N'\le_\lambda N$ means that $N'$ is represented by some
$\lambda'\le\lambda$.
\end{lemma}

\begin{proof}
Set $A_s=S_s^{-1}=(1+s^2\Delta)^3$ and $U=S_sW$.  Since
$\Pi A_sU=\Pi W=0$, the distributional identity
\[
 \Pi U=S_s[A_s,\Pi]U
\]
holds.  The commutator expands as a sum of
$\binom3j s^{2j}[\Delta^j,\Pi]$, $1\le j\le3$.
Because the principal symbol of $\Delta$ is scalar,
$[\Delta^j,\Pi]$ is a differential operator of order at most $2j-1$.
The spectral theorem and elliptic Sobolev equivalence give
\[
 \|S_s\|_{L^2\to H^m}\le C_ms^{-m},\qquad 0\le m\le6.
\]
If $B_j=[\Delta^j,\Pi]$, its formal adjoint has the same order, so
\[
 \|S_s s^{2j}B_j\|_{L^2\to L^2}
 =\|s^{2j}B_j^*S_s\|_{L^2\to L^2}
 \le Cs^{2j}s^{-(2j-1)}=Cs.
\]
This proves \Cref{eq:current-projection-leakage}, including its meaning
for distributional $W$ by the bounded extension of the displayed
operator.

Apply this result to $R N'=0$ and $(1-R)Z=0$.  With $u=S_sN'$ and
$v=S_sZ$, the orthogonal splitting gives
\[
 |\langle u,v\rangle_2|
 \le\|Ru\|_2\|Rv\|_2+
       \|(1-R)u\|_2\|(1-R)v\|_2
 \le Cs\|u\|_2\|v\|_2.
\]
Now use \Cref{eq:current-submeasure-ltwo}.
\end{proof}

The topological pairing used next applies to the distributional
coefficients just constructed.  If $C_1,C_2\in H^{-3}$ are closed
currents, Hodge decomposition of the closed distributional forms $*C_i$
and commutation of $S_s$ with $\dd$ give
\begin{equation}\label{eq:current-topological-pairing}
 \langle S_sC_1,*S_sC_2\rangle_2
  =\int_X(*S_sC_1)\wedge(*S_sC_2)
  =[*C_1]\cdot[*C_2].
\end{equation}
For the underlying Green-operator decomposition, see
\cite[Sections~1--2, equations~(1.20)--(1.22) and~(2.1)--(2.2)]{TaylorHodge2010}.
Its extension to distributions follows from the Sobolev mapping
properties of the Green operator. To see the regularity involved,
write $*C_i=h_i+\dd a_i$, with $h_i$
harmonic and $a_i\in H^{-2}$.  Then $S_sh_i=h_i$ and
$S_sa_i\in H^4$.  Integration by parts eliminates all terms containing
an exact factor, leaving $\int_Xh_1\wedge h_2$.  Thus no product of the
unsmoothed distributions is used in
\Cref{eq:current-topological-pairing}.

\begin{proposition}[Angular energy]\label{prop:angular-energy}
In setup \textup{(A)} or \textup{(B)}, the correction satisfies
$Q\in L^2(X,E)$ and
\begin{equation}\label{eq:current-angular-energy}
 \iint_{\dist(y,z)<s}|\ell-\tau_{z\to y}m|^2
           \,\dd\lambda(y,\ell)\dd\lambda(z,m)\le Cs^4.
\end{equation}
The estimates in
\Cref{lem:current-positive-pairing,lem:current-commutator} hold for all
the indicated submeasures, and, for every such fixed submeasure current
$N'$,
\begin{equation}\label{eq:current-mixed-vanishing}
 \lim_{s\downarrow0}\langle S_sN',S_sQ\rangle_2=0.
\end{equation}
In fact the convergence in \Cref{eq:current-mixed-vanishing} is uniform
over $\lambda'\le\lambda$.
\end{proposition}

\begin{proof}
Let $\mathcal A_s$ denote the double integral in
\Cref{eq:current-angular-energy} and put $x_s=\|S_sQ\|_2$.
Since $Q$ is self-dual and $S_s$ commutes with $*$,
\Cref{eq:current-topological-pairing} applied to $T=N+Q$ gives
\[
 [*T]^2
  =\langle S_sN,*S_sN\rangle_2
       +2\langle S_sN,S_sQ\rangle_2+x_s^2.
\]
Use \Cref{eq:current-positive-pairing,eq:current-bracket-bound} for
the first term and \Cref{eq:current-mixed-bound} for the second.  The
result is
\begin{equation}\label{eq:current-energy-coercivity}
 cs^{-4}\mathcal A_s+x_s^2-Cx_s\le C+[*T]^2.
\end{equation}
The right side is independent of $s$.  In setup \textup{(A)} its
intersection term is zero.  In setup \textup{(B)} that term is bounded
uniformly for mass-one $N$: its harmonic coefficients are pairings with
fixed smooth forms and are bounded by \Cref{eq:current-line-measure}.
Completing the square in $x_s$ bounds both $x_s$ and
$s^{-4}\mathcal A_s$.

As $s\downarrow0$, $S_sQ$ converges to $Q$ distributionally.  A weakly
convergent subsequence of the bounded $L^2$ family therefore has limit
$Q$, proving $Q\in L^2$.  Its anti-invariance is preserved in this
limit, and the spectral theorem now gives $S_sQ\to Q$ strongly in $L^2$.

Choose smooth anti-invariant $Q_j$ with $\|Q_j-Q\|_2\to0$; for example
apply $R$ to a smooth approximation.  By
\Cref{eq:current-mixed-bound} and the $L^2$ contraction property of
$S_s$,
\[
 |\langle S_sN',S_s(Q-Q_j)\rangle_2|
       \le C\|Q-Q_j\|_2.
\]
For fixed $j$, $S_s^2Q_j\to Q_j$ in $C^\infty$, whereas
$N'(Q_j)=0$.  Therefore
\[
 \langle S_sN',S_sQ_j\rangle_2
       =N'(S_s^2Q_j-Q_j)\longrightarrow0.
\]
The last convergence is uniform over the submeasures, because their
trace masses are at most one.  First let $s$ tend to zero and then $j$
to infinity to obtain \Cref{eq:current-mixed-vanishing}.
\end{proof}
